% Generated draft: requires mathematical and source review.
\documentclass[11pt]{article}
\usepackage[margin=25mm]{geometry}
\usepackage{fontspec}
\setmainfont{Latin Modern Roman}
\usepackage{amsmath,amssymb,mathtools,unicode-math}
\setmathfont{Latin Modern Math}
\usepackage{enumitem,needspace,xurl,hyperref}
\setlength{\parindent}{0pt}
\setlength{\parskip}{5pt}
\setlength{\emergencystretch}{3em}
\newcommand{\N}{\mathbb{N}}
\newcommand{\Z}{\mathbb{Z}}
\newcommand{\Q}{\mathbb{Q}}
\newcommand{\R}{\mathbb{R}}
\newcommand{\C}{\mathbb{C}}
\newcommand{\F}{\mathbb{F}}
\newcommand{\A}{\mathbb{A}}
\newcommand{\PP}{\mathbb{P}}
\newcommand{\E}{\mathbb{E}}
\newcommand{\eps}{\varepsilon}
\newcommand{\dd}{\,\mathrm{d}}
\begin{document}
\section*{UnsolvedMath ID: 20001963}
\textbf{Source problem code:} catalogue entry
\subsection{Definitions and mathematical statement}
Hausel's Conjecture 6.1 asks for surjectivity of the hyperk\"ahler Kirwan map
\[
\kappa:H_G^*(T^*A)\cong H^*(BG)\longrightarrow H^*(T^*A\mathbin{///}_{\xi}G).
\]
In a smooth finite-dimensional version, \(G\) is compact, \(T^*A\) has a tri-Hamiltonian action with moment map \(\mu\), \(\xi\) is an invariant regular value, and \(G\) acts freely on \(\mu^{-1}(\xi)\); the quotient is \(\mu^{-1}(\xi)/G\). Restriction followed by the free-action identification defines \(\kappa\). The displayed \(H^*(BG)\) identification requires contractibility of \(T^*A\), as in a linear or affine model. The extract omits the definition of \(A\), coefficient ring, and precise regularity, isotropy and possibly gauge-theoretic hypotheses, so this version is explicit rather than silently universal.

The supplied aid treats an integral special case. For integers \(1\leq k\leq n\), take \(A=\operatorname{Hom}(\C^k,\C^n)\), pairs \((I,J)\) in \(T^*A\), and the action \((I,J)\mapsto(Ig^{-1},gJ)\) of \(U(k)\). Choose moment-map sign and scale so the positive real level is
\[
I^*I-JJ^*=\zeta\operatorname{Id}_k,\qquad JI=0,\qquad\zeta>0.
\]
The aid identifies its quotient with \(T^*\operatorname{Gr}(k,n)\). Since this retracts onto the Grassmannian, the integral Kirwan map is, with the tautological-subbundle convention,
\[
\Z[c_1,\ldots,c_k]\longrightarrow H^*(\operatorname{Gr}(k,n);\Z),
\qquad c_i\longmapsto c_i(S),\quad |c_i|=2i.
\]
Define \(b_r\) by
\[
\sum_{r\geq0}b_rt^r=(1+c_1t+\cdots+c_kt^k)^{-1}.
\]
The reported kernel is \((b_{n-k+1},\ldots,b_n)\). Schubert Giambelli polynomials give an additive section: for a partition \(\lambda\) inside the \(k\times(n-k)\) rectangle, lift its Schubert class to \(\det(b_{\lambda_i-i+j})_{1\leq i,j\leq k}\), with \(b_0=1\) and \(b_r=0\) for \(r<0\), using the corresponding Schubert convention. This is an additive, not a ring, splitting. The first nonzero kernel occurs in cohomological degree \(2(n-k+1)\), where it is the rank-one group generated by \(b_{n-k+1}\). This special-case presentation does not prove the full conjecture.

The historical Comment 6.2 reports results for \(M^1_{\mathrm{Dol}}(\mathrm{GL}(2,\C))\) and \(M^d_{\mathrm{Dol}}(\mathrm{GL}(n,\C))\), a quiver-variety conjecture, and proved hyperpolygon and hypertoric cases, with the references below. Their stability, degree and smoothness hypotheses are not present in the extract, so the historical list is not a blanket theorem for arbitrary parameters.

The appended integration section calls an oriented circle manifold circle-compact when \(M^{U(1)}\) is compact, and defines
\[
\widehat H_{U(1)}^*(M)=H_{U(1)}^*(M;\Q)\otimes_{\Q[u]}\Q(u),\qquad |u|=2,
\]
\[
\int_M\alpha:=\sum_{F\subset M^{U(1)}}\int_F
\frac{i_F^*\alpha}{e_{U(1)}(N_F)}\in\Q(u).
\]
The sum is over fixed components, with compatible orientations; nonzero normal circle weights make the Euler classes invertible after localization. This is a localized definition, not ordinary integration of arbitrary forms over noncompact \(M\). The source says fiber rotation induces circle-compact actions on its examples; descent at a general moment-map value needs checking, since fiber rotation scales the complex moment map. Wu mentions Prato--Wu's tempered-distribution interpretation, and Hausel asks for a formulation using equivariant differential forms. That additional question is not answered by the Grassmannian calculation.
\subsection{Short English statement}
The source conjectures hyperk\"ahler Kirwan surjectivity. The supplied aid gives an integral cotangent-Grassmannian presentation, its exact inverse-Chern kernel and an additive Schubert splitting. The appended localized integration definition and differential-form question are separate and retain their coefficient and compact-fixed-set qualifications.
\subsection{Sources}
\begin{itemize}
\item T. Hausel, Conjecture 6.1 and Comment 6.6; T. Hausel and N. Proudfoot, Definition 6.4; S. Wu, Comment 6.5, in the AIM workshop \emph{Moment maps and surjectivity in various geometries}: \url{https://aimath.org/WWN/momentmaps/momentmaps.pdf}.
\item Historical references as supplied: Hausel--Thaddeus (2000), Markman (2001), Nakajima (2002), Konno (2000), Hausel--Proudfoot (2003), Hausel--Sturmfels (2002); E. Prato and S. Wu [PW94]. Full entries and full theorem hypotheses are absent.
\item Local registry \texttt{data/unsolvedmath/problems.json}, permanent ID 20001963, code \texttt{catalogue entry}; canonical \texttt{aim-geometry-notes.json}, \texttt{notes[300]}, source item 6.1; updated 14 August 2026.
\item Attempt 1: partial result, partially solved status, rigor label proved, novelty confidence low. The aid is machine generated, not peer reviewed, and requires independent verification before citation.
\end{itemize}
\end{document}
